\documentclass[aspectratio=169,10pt]{beamer}
\makeatletter\def\input@path{{../../estilo/}{../estilo/}{estilo/}}\makeatother
\usepackage{devmobile}
\modulo[Kotlin]{05}{Interfaces com Jetpack Compose}{UI declarativa, componentes, layouts, modifiers, listas, Material 3 e temas}

\begin{document}
\capa

\begin{frame}{Objetivos do módulo}
  \begin{conceito}[Ao final deste módulo você será capaz de]
    \begin{itemize}
      \item Explicar o paradigma declarativo, o layout por restrições e a recomposição.
      \item Usar os componentes básicos: textos, botões, imagens, campos e cartões.
      \item Montar telas com \cod{Column}, \cod{Row} e \cod{Box}, alinhamentos e espaçamentos.
      \item Aplicar \cod{Modifier} entendendo que a ordem importa.
      \item Exibir listas eficientes com \cod{LazyColumn} e grades.
      \item Estruturar telas com \cod{Scaffold} e personalizar o tema Material 3.
    \end{itemize}
  \end{conceito}
\end{frame}

\sumario

% =============================================================================
%  Fundamentos compartilhados --- Módulo 05: interfaces declarativas
%  Incluído pelas trilhas Kotlin e Flutter com % =============================================================================
%  Fundamentos compartilhados --- Módulo 05: interfaces declarativas
%  Incluído pelas trilhas Kotlin e Flutter com % =============================================================================
%  Fundamentos compartilhados --- Módulo 05: interfaces declarativas
%  Incluído pelas trilhas Kotlin e Flutter com \input{../fundamentos/f05-ui-declarativa}
% =============================================================================
\section{Fundamentos: interfaces declarativas}

\begin{frame}{Interfaces imperativas e declarativas}
  \begin{columns}[T]
    \begin{column}{0.49\textwidth}
      \begin{definicao}[Interface imperativa]
        O programa cria os componentes uma vez e depois os \textbf{altera} passo a passo (\cod{setText}, \cod{setVisibility}). Basta esquecer um deles para a tela mostrar dados errados.
      \end{definicao}
    \end{column}
    \begin{column}{0.49\textwidth}
      \begin{definicao}[Interface declarativa]
        O programa \textbf{descreve} como a tela deve ser para cada estado. Quando o estado muda, o framework gera a nova descrição e aplica apenas as diferenças.
      \end{definicao}
    \end{column}
  \end{columns}
  \vspace{2pt}
  \centering
  \begin{tikzpicture}[node distance=2mm and 6mm,every node/.style={font=\scriptsize}]
    \node[rotulo,anchor=east,text width=1.6cm,align=right] (li) at (0,0) {\textbf{imperativo}};
    \node[caixa,right=3mm of li,minimum width=1.6cm] (e1) {estado muda};
    \node[caixaambar,right=of e1,minimum width=3.2cm] (c1) {seu código acha e altera\\ cada componente afetado};
    \node[caixa,right=of c1,minimum width=2.2cm] (t1) {tela (mutável)};
    \draw[seta] (e1) -- (c1); \draw[seta] (c1) -- (t1);
    \node[rotulo,anchor=east,text width=1.6cm,align=right] (ld) at (0,-1.05) {\textbf{declarativo}};
    \node[caixa,right=3mm of ld,minimum width=1.6cm] (e2) {estado muda};
    \node[caixapetroleo,right=of e2,minimum width=2.1cm] (f2) {UI = f(estado)};
    \node[caixa,right=of f2,minimum width=2.1cm] (d2) {nova descrição};
    \node[caixapetroleo,right=of d2,minimum width=2.6cm] (a2) {framework compara\\ e aplica diferenças};
    \draw[seta] (e2) -- (f2); \draw[seta] (f2) -- (d2); \draw[seta] (d2) -- (a2);
  \end{tikzpicture}
  \raggedright
  \begin{historia}
    Views do Android e UIKit do iOS (2008) são imperativos. O React (2013) popularizou a UI declarativa, depois adotada pelo Flutter (2018), SwiftUI (2019) e Jetpack Compose (2021).
  \end{historia}
\end{frame}

\begin{frame}{Composição: a tela como uma árvore}
  \begin{columns}[c]
    \begin{column}{0.52\textwidth}
      \centering
      \begin{tikzpicture}[level distance=10mm,
          every node/.style={caixa,font=\tiny,minimum height=5mm,inner sep=2pt},
          edge from parent/.style={draw=azulpetroleo,line width=0.6pt},
          level 1/.style={sibling distance=30mm},
          level 2/.style={sibling distance=21mm},
          level 3/.style={sibling distance=10mm}]
        \node[caixaescura,font=\tiny] (raiz) {TelaEquipamentos}
          child {node {BarraSuperior}}
          child {node {ListaEquipamentos}
            child {node[caixaambar,font=\tiny] {CartaoBancada}
              child {node {Ícone}}
              child {node {Textos}}
              child {node {Barra}}}
            child {node[caixaambar,font=\tiny] {CartaoBancada}}};
        \draw[seta] (-2.6,0.2) -- (-2.6,-3.2);
        \draw[setaambar] (-2.2,-3.2) -- (-2.2,0.2);
        \node[rotulo,anchor=north west,text width=6cm,align=left] at (-2.75,-3.45) {\textcolor{azulpetroleo}{$\downarrow$} dados descem por parâmetros\quad \textcolor{ambar}{$\uparrow$} eventos sobem por \emph{callbacks}};
      \end{tikzpicture}
    \end{column}
    \begin{column}{0.46\textwidth}
      \begin{definicao}[Composição]
        Construir interfaces complexas \textbf{combinando} componentes pequenos e independentes, em vez de herdar de componentes grandes. Cada componente recebe dados e devolve uma descrição de UI.
      \end{definicao}
      \begin{analogia}
        Como peças de montar: o cartão é feito de ícone, textos e indicador; a lista, de vários cartões. Trocar uma peça não obriga a refazer as outras.
      \end{analogia}
      \begin{dica}[Um bom componente]
        Faz uma coisa só, recebe tudo por parâmetro, não sabe de onde vêm os dados e pode ser visto isoladamente em uma \emph{preview}.
      \end{dica}
    \end{column}
  \end{columns}
\end{frame}

\begin{frame}{Layout por restrições}
  \begin{columns}[c]
    \begin{column}{0.47\textwidth}
      \centering
      \begin{tikzpicture}[node distance=11mm]
        \node[caixaescura,minimum width=3.4cm] (t) {Tela \quad 360 $\times$ 800};
        \node[caixa,minimum width=3.4cm,below=of t] (c) {Coluna};
        \node[caixaambar,minimum width=3.4cm,below=of c] (k) {Cartão};
        \draw[seta] ([xshift=-10mm]t.south) -- node[rotulo,left,text width=2.1cm,align=right]{largura 0--360\\ altura 0--800} ([xshift=-10mm]c.north);
        \draw[seta] ([xshift=-10mm]c.south) -- node[rotulo,left,text width=2.1cm,align=right]{largura 0--360\\ altura livre ($\infty$)} ([xshift=-10mm]k.north);
        \draw[setaambar] ([xshift=10mm]k.north) -- node[rotulo,right,text width=1.6cm,align=left]{``tenho\\ 360 $\times$ 96''} ([xshift=10mm]c.south);
        \node[rotulo,below=1.5mm of k,text width=4cm] {a coluna põe o cartão em $y = 0$\\ e o próximo em $y = 96$};
      \end{tikzpicture}
    \end{column}
    \begin{column}{0.51\textwidth}
      \begin{definicao}[Restrições (\emph{constraints})]
        Intervalo de tamanhos que o pai aceita para o filho: largura e altura mínimas e máximas. \textbf{Justa} quando mínimo = máximo; \textbf{frouxa} quando o mínimo é zero.
      \end{definicao}
      \begin{conceito}[O algoritmo em três frases]
        \textbf{Restrições descem}: o pai diz os limites. \textbf{Tamanhos sobem}: o filho escolhe um tamanho dentro deles. \textbf{O pai posiciona}: o filho não sabe onde está.
      \end{conceito}
      \begin{pordentro}[Uma passada só]
        Cada nó é medido uma única vez por quadro, o que mantém o layout rápido mesmo com milhares de nós. No Compose, cada \cod{Modifier} ajusta as restrições antes de repassá-las; no Flutter, \cod{Padding} e \cod{SizedBox} fazem esse papel.
      \end{pordentro}
    \end{column}
  \end{columns}
\end{frame}

\begin{frame}{Do estado aos pixels: o caminho de um quadro}
  \centering
  \begin{tikzpicture}[node distance=5mm,every node/.style={font=\scriptsize}]
    \node[caixaambar,minimum width=2cm,minimum height=11mm] (s) {estado\\ muda};
    \node[caixapetroleo,right=of s,minimum width=2.3cm,minimum height=11mm] (c) {\textbf{composição}\\ \tiny o que mostrar};
    \node[caixapetroleo,right=of c,minimum width=2.3cm,minimum height=11mm] (l) {\textbf{layout}\\ \tiny tamanho e posição};
    \node[caixapetroleo,right=of l,minimum width=2.3cm,minimum height=11mm] (d) {\textbf{desenho}\\ \tiny cores, textos, formas};
    \node[caixaescura,right=of d,minimum width=2.3cm,minimum height=11mm] (r) {\textbf{rasterização}\\ \tiny GPU gera os pixels};
    \draw[seta] (s) -- (c); \draw[seta] (c) -- (l); \draw[seta] (l) -- (d); \draw[seta] (d) -- (r);
    \node[rotulo,below=1mm of c] {Compose: \emph{Composition}\\ Flutter: \emph{build}};
    \node[rotulo,below=1mm of l] {\emph{Layout}\\ \emph{layout}};
    \node[rotulo,below=1mm of d] {\emph{Drawing}\\ \emph{paint} e \emph{composite}};
    \node[rotulo,below=1mm of r] {motor gráfico\\ (Skia, Impeller)};
  \end{tikzpicture}
  \raggedright
  \begin{columns}[T]
    \begin{column}{0.49\textwidth}
      \begin{definicao}[Quadro e orçamento de tempo]
        A tela é redesenhada no ritmo do sinal vertical do display (VSync): a 60\,Hz cada quadro tem \textbf{16,7\,ms}; a 120\,Hz, só \textbf{8,3\,ms}. Se o trabalho não couber, o quadro se perde e a animação engasga (\emph{jank}).
      \end{definicao}
    \end{column}
    \begin{column}{0.49\textwidth}
      \begin{pordentro}[Pular fases economiza tempo]
        Uma cor que anima só precisa de desenho; uma posição que muda, de layout e desenho. Ler o estado na fase mais tardia possível evita repetir as anteriores.
      \end{pordentro}
      \begin{dica}
        O \emph{Performance} do Flutter DevTools e o \emph{Layout Inspector} do Android Studio mostram onde o tempo do quadro foi gasto.
      \end{dica}
    \end{column}
  \end{columns}
\end{frame}

\begin{frame}{Por dentro: descrição, instância e renderização}
  \begin{columns}[T]
    \begin{column}{0.56\textwidth}
      \scriptsize
      \begin{tabularx}{\linewidth}{@{}L l l@{}}
        \toprule
        \cabecalho{Papel} & \cabecalho{Flutter} & \cabecalho{Compose} \\
        \midrule
        Descrição leve, recriada a cada mudança & \cod{Widget} & chamada \cod{@Composable} \\
        Instância duradoura, guarda o estado & \cod{Element} + \cod{State} & composição (\cod{remember}) \\
        Mede, posiciona e desenha & \cod{RenderObject} & \cod{LayoutNode} \\
        Informa leitores de tela & árvore de semântica & árvore de semântica \\
        \bottomrule
      \end{tabularx}
      \normalsize
      \begin{analogia}
        O widget é a planta do arquiteto: barata de redesenhar. O elemento é o apartamento construído, onde mora o estado. O objeto de renderização é a estrutura medida e pintada. Mudar a planta não obriga a demolir o prédio.
      \end{analogia}
    \end{column}
    \begin{column}{0.42\textwidth}
      \begin{conceito}[Por que separar]
        Descrições são criadas aos milhares por segundo, então precisam ser baratas. Medir e desenhar é caro, então os objetos de renderização são \textbf{reaproveitados} enquanto possível.
      \end{conceito}
      \begin{pordentro}[Semântica e acessibilidade]
        Em paralelo, o framework monta uma árvore de semântica (rótulos, papéis, ações) que o TalkBack e o VoiceOver leem. Ícone sem rótulo fica mudo para quem não enxerga (Módulo 04).
      \end{pordentro}
    \end{column}
  \end{columns}
\end{frame}

\begin{frame}{Atualizar sem refazer tudo}
  \begin{columns}[c]
    \begin{column}{0.42\textwidth}
      \centering
      \begin{tikzpicture}[level distance=10mm,
          every node/.style={caixa,font=\tiny,minimum height=5mm,minimum width=9mm,inner sep=2pt},
          edge from parent/.style={draw=azulpetroleo,line width=0.6pt},
          level 1/.style={sibling distance=22mm},
          level 2/.style={sibling distance=11mm}]
        \node {Tela}
          child {node {Barra} child {node {Título}}}
          child {node[caixaambarforte] (m) {Medidor}
            child {node[caixaambar] {Valor}}
            child {node[caixaambar] {Barra}}};
        \node[rotulo,below=11mm of m,text width=3.4cm] {o estado \cod{temperatura} mudou: só o \textbf{Medidor} e seus filhos são reconstruídos};
      \end{tikzpicture}
    \end{column}
    \begin{column}{0.56\textwidth}
      \begin{conceito}[Flutter: \emph{rebuild} e reconciliação]
        \cod{setState} marca o elemento como sujo; no próximo quadro, só ele executa \cod{build}. Filhos de mesmo tipo e chave são reaproveitados e só atualizados.
      \end{conceito}
      \begin{conceito}[Compose: recomposição inteligente]
        O Compose sabe quais funções leram cada estado e só executa essas de novo; as que recebem parâmetros iguais aos anteriores são puladas (\emph{skipping}).
      \end{conceito}
      \begin{atencao}[Consequência importante]
        Funções de interface podem rodar muitas vezes, em qualquer ordem: devem ser \textbf{rápidas e sem efeitos colaterais}. Rede e banco ficam em APIs de efeito (Módulo 07).
      \end{atencao}
    \end{column}
  \end{columns}
\end{frame}

\begin{frame}[fragile]{Identidade e chaves}
  \begin{columns}[T]
    \begin{column}{0.47\textwidth}
      \centering
      \begin{tikzpicture}[baseline=(current bounding box.north),every node/.style={font=\tiny}]
        \tikzset{item/.style={caixa,minimum width=1.5cm,minimum height=5mm,inner sep=1pt}}
        \node[rotulo] at (0.75,0.55) {antes};
        \node[item,caixaambar] (a) at (0.75,0) {A \quad expandido};
        \node[item] at (0.75,-0.6) {B};
        \node[item] at (0.75,-1.2) {C};
        \node[rotulo,text width=2cm] at (2.95,0.55) {remove A, \textbf{sem} chave};
        \node[item,caixaambar,draw=vermelhoerro] at (2.95,0) {B \quad expandido};
        \node[item] at (2.95,-0.6) {C};
        \node[rotulo,text width=2cm] at (5.15,0.55) {remove A, \textbf{com} chave};
        \node[item] at (5.15,0) {B};
        \node[item] at (5.15,-0.6) {C};
        \node[rotulo,text width=2cm,text=vermelhoerro] at (2.95,-1.25) {o estado ficou na posição 1};
        \node[rotulo,text width=2cm,text=verdeok] at (5.15,-1.25) {o estado saiu junto com A};
      \end{tikzpicture}
\begin{codigo}[title=Compose]{Kotlin}
items(itens, key = { it.id }) { eq ->
  ItemEquipamento(eq)
}
\end{codigo}
\begin{codigo}[title=Flutter]{Dart}
ItemEquipamento(
    key: ValueKey(eq.id), eq: eq, onTap: abrir)
\end{codigo}
    \end{column}
    \begin{column}{0.51\textwidth}
      \begin{definicao}[Chave (\emph{key})]
        Identificador estável que diz ao framework qual item é qual entre duas versões da árvore. Sem chave, a correspondência é feita pela \textbf{posição}.
      \end{definicao}
      \begin{conceito}[Quando importa]
        Em listas cujos itens têm estado próprio (expandido, campo digitado, animação) e podem ser inseridos, removidos ou reordenados.
      \end{conceito}
      \begin{dica}
        Use um identificador do domínio, como o \cod{id} do equipamento. Nunca o índice da lista nem um valor aleatório criado a cada reconstrução.
      \end{dica}
    \end{column}
  \end{columns}
\end{frame}

\begin{frame}{Temas, \emph{tokens} e contexto implícito}
  \begin{columns}[c]
    \begin{column}{0.42\textwidth}
      \centering
      \begin{tikzpicture}[node distance=7mm and 3mm,every node/.style={font=\tiny}]
        \node[caixaambarforte,minimum width=2.4cm,font=\tiny\bfseries] (t) {Tema};
        \node[caixa,below=of t,minimum width=1.6cm,font=\tiny] (s) {Tela};
        \node[caixa,below left=of s,minimum width=1.3cm,font=\tiny] (l) {Lista};
        \node[caixa,below right=of s,minimum width=1.3cm,font=\tiny] (b) {Botão};
        \node[caixa,below=of l,minimum width=1.3cm,font=\tiny] (c) {Cartão};
        \draw[seta] (t) -- (s); \draw[seta] (s) -- (l); \draw[seta] (s) -- (b); \draw[seta] (l) -- (c);
        \draw[setaambar,dashed] (t.east) -| (b.north);
        \draw[setaambar,dashed] (t.west) -| ([xshift=-3mm]c.west) -- (c.west);
        \node[rotulo,text width=4cm,below=2mm of c.south east] {tracejado: componentes leem o tema direto do contexto, sem receber parâmetros};
      \end{tikzpicture}
    \end{column}
    \begin{column}{0.56\textwidth}
      \begin{definicao}[\emph{Tokens} de design]
        Decisões visuais nomeadas pelo \textbf{papel}, não pelo valor: \cod{primary}, \cod{surface}, \cod{titleMedium}. O componente pede ``a cor primária''; o tema decide qual é.
      \end{definicao}
      \begin{pordentro}[Contexto implícito]
        Tema, idioma e tamanho da tela ficam disponíveis para toda uma subárvore: \cod{InheritedWidget} no Flutter (\cod{Theme.of(context)}), \cod{CompositionLocal} no Compose. Quem lê é atualizado quando o valor muda.
      \end{pordentro}
      \begin{dica}
        Com \emph{tokens}, tema escuro e cor dinâmica (Android 12+) funcionam sem mexer nos componentes.
      \end{dica}
    \end{column}
  \end{columns}
\end{frame}

\begin{frame}{Pixels, densidade e telas de todos os tamanhos}
  \begin{columns}[T]
    \begin{column}{0.49\textwidth}
      \centering
      \begin{tikzpicture}[baseline=(current bounding box.north)]
        \draw[cinzamedio,step=0.2] (0,0) grid (1.2,1.2);
        \fill[ambar,opacity=0.7] (0.2,0.2) rectangle (1.0,1.0);
        \node[rotulo] at (0.6,-0.25) {160 dpi: 48\,dp = 48\,px};
        \draw[cinzamedio,step=0.0667] (2.6,0) grid (3.8,1.2);
        \fill[ambar,opacity=0.7] (2.8,0.2) rectangle (3.6,1.0);
        \node[rotulo] at (3.2,-0.25) {480 dpi: 48\,dp = 144\,px};
      \end{tikzpicture}
      \raggedright
      \begin{definicao}[Pixel independente de densidade]
        \cod{px = dp} $\times$ \cod{(dpi} $\div$ \cod{160)}: um botão de 48\,dp tem o mesmo tamanho físico em qualquer tela. No Flutter, é o pixel lógico (fator em \cod{devicePixelRatio}).
      \end{definicao}
      \begin{dica}[Texto]
        Fontes também seguem a escala escolhida pelo usuário: \cod{sp} no Android, \cod{TextScaler} no Flutter.
      \end{dica}
    \end{column}
    \begin{column}{0.49\textwidth}
      \begin{conceito}[Responsivo $\times$ adaptativo]
        \textbf{Responsivo}: o mesmo layout se estica ou reorganiza. \textbf{Adaptativo}: um layout por classe de tamanho de janela (barra inferior no celular, trilho lateral no tablet). Pense em \textbf{espaço disponível}, não em modelo de aparelho.
      \end{conceito}
      \begin{atencao}[Android 16 e telas grandes]
        Com \cod{targetSdk} 36 (exigido pela Play desde agosto de 2026), telas com largura mínima $\geq$ 600\,dp ignoram travas de orientação e de redimensionamento; no API 37 não haverá exceção.
      \end{atencao}
    \end{column}
  \end{columns}
\end{frame}

% =============================================================================
\section{Fundamentos: interfaces declarativas}

\begin{frame}{Interfaces imperativas e declarativas}
  \begin{columns}[T]
    \begin{column}{0.49\textwidth}
      \begin{definicao}[Interface imperativa]
        O programa cria os componentes uma vez e depois os \textbf{altera} passo a passo (\cod{setText}, \cod{setVisibility}). Basta esquecer um deles para a tela mostrar dados errados.
      \end{definicao}
    \end{column}
    \begin{column}{0.49\textwidth}
      \begin{definicao}[Interface declarativa]
        O programa \textbf{descreve} como a tela deve ser para cada estado. Quando o estado muda, o framework gera a nova descrição e aplica apenas as diferenças.
      \end{definicao}
    \end{column}
  \end{columns}
  \vspace{2pt}
  \centering
  \begin{tikzpicture}[node distance=2mm and 6mm,every node/.style={font=\scriptsize}]
    \node[rotulo,anchor=east,text width=1.6cm,align=right] (li) at (0,0) {\textbf{imperativo}};
    \node[caixa,right=3mm of li,minimum width=1.6cm] (e1) {estado muda};
    \node[caixaambar,right=of e1,minimum width=3.2cm] (c1) {seu código acha e altera\\ cada componente afetado};
    \node[caixa,right=of c1,minimum width=2.2cm] (t1) {tela (mutável)};
    \draw[seta] (e1) -- (c1); \draw[seta] (c1) -- (t1);
    \node[rotulo,anchor=east,text width=1.6cm,align=right] (ld) at (0,-1.05) {\textbf{declarativo}};
    \node[caixa,right=3mm of ld,minimum width=1.6cm] (e2) {estado muda};
    \node[caixapetroleo,right=of e2,minimum width=2.1cm] (f2) {UI = f(estado)};
    \node[caixa,right=of f2,minimum width=2.1cm] (d2) {nova descrição};
    \node[caixapetroleo,right=of d2,minimum width=2.6cm] (a2) {framework compara\\ e aplica diferenças};
    \draw[seta] (e2) -- (f2); \draw[seta] (f2) -- (d2); \draw[seta] (d2) -- (a2);
  \end{tikzpicture}
  \raggedright
  \begin{historia}
    Views do Android e UIKit do iOS (2008) são imperativos. O React (2013) popularizou a UI declarativa, depois adotada pelo Flutter (2018), SwiftUI (2019) e Jetpack Compose (2021).
  \end{historia}
\end{frame}

\begin{frame}{Composição: a tela como uma árvore}
  \begin{columns}[c]
    \begin{column}{0.52\textwidth}
      \centering
      \begin{tikzpicture}[level distance=10mm,
          every node/.style={caixa,font=\tiny,minimum height=5mm,inner sep=2pt},
          edge from parent/.style={draw=azulpetroleo,line width=0.6pt},
          level 1/.style={sibling distance=30mm},
          level 2/.style={sibling distance=21mm},
          level 3/.style={sibling distance=10mm}]
        \node[caixaescura,font=\tiny] (raiz) {TelaEquipamentos}
          child {node {BarraSuperior}}
          child {node {ListaEquipamentos}
            child {node[caixaambar,font=\tiny] {CartaoBancada}
              child {node {Ícone}}
              child {node {Textos}}
              child {node {Barra}}}
            child {node[caixaambar,font=\tiny] {CartaoBancada}}};
        \draw[seta] (-2.6,0.2) -- (-2.6,-3.2);
        \draw[setaambar] (-2.2,-3.2) -- (-2.2,0.2);
        \node[rotulo,anchor=north west,text width=6cm,align=left] at (-2.75,-3.45) {\textcolor{azulpetroleo}{$\downarrow$} dados descem por parâmetros\quad \textcolor{ambar}{$\uparrow$} eventos sobem por \emph{callbacks}};
      \end{tikzpicture}
    \end{column}
    \begin{column}{0.46\textwidth}
      \begin{definicao}[Composição]
        Construir interfaces complexas \textbf{combinando} componentes pequenos e independentes, em vez de herdar de componentes grandes. Cada componente recebe dados e devolve uma descrição de UI.
      \end{definicao}
      \begin{analogia}
        Como peças de montar: o cartão é feito de ícone, textos e indicador; a lista, de vários cartões. Trocar uma peça não obriga a refazer as outras.
      \end{analogia}
      \begin{dica}[Um bom componente]
        Faz uma coisa só, recebe tudo por parâmetro, não sabe de onde vêm os dados e pode ser visto isoladamente em uma \emph{preview}.
      \end{dica}
    \end{column}
  \end{columns}
\end{frame}

\begin{frame}{Layout por restrições}
  \begin{columns}[c]
    \begin{column}{0.47\textwidth}
      \centering
      \begin{tikzpicture}[node distance=11mm]
        \node[caixaescura,minimum width=3.4cm] (t) {Tela \quad 360 $\times$ 800};
        \node[caixa,minimum width=3.4cm,below=of t] (c) {Coluna};
        \node[caixaambar,minimum width=3.4cm,below=of c] (k) {Cartão};
        \draw[seta] ([xshift=-10mm]t.south) -- node[rotulo,left,text width=2.1cm,align=right]{largura 0--360\\ altura 0--800} ([xshift=-10mm]c.north);
        \draw[seta] ([xshift=-10mm]c.south) -- node[rotulo,left,text width=2.1cm,align=right]{largura 0--360\\ altura livre ($\infty$)} ([xshift=-10mm]k.north);
        \draw[setaambar] ([xshift=10mm]k.north) -- node[rotulo,right,text width=1.6cm,align=left]{``tenho\\ 360 $\times$ 96''} ([xshift=10mm]c.south);
        \node[rotulo,below=1.5mm of k,text width=4cm] {a coluna põe o cartão em $y = 0$\\ e o próximo em $y = 96$};
      \end{tikzpicture}
    \end{column}
    \begin{column}{0.51\textwidth}
      \begin{definicao}[Restrições (\emph{constraints})]
        Intervalo de tamanhos que o pai aceita para o filho: largura e altura mínimas e máximas. \textbf{Justa} quando mínimo = máximo; \textbf{frouxa} quando o mínimo é zero.
      \end{definicao}
      \begin{conceito}[O algoritmo em três frases]
        \textbf{Restrições descem}: o pai diz os limites. \textbf{Tamanhos sobem}: o filho escolhe um tamanho dentro deles. \textbf{O pai posiciona}: o filho não sabe onde está.
      \end{conceito}
      \begin{pordentro}[Uma passada só]
        Cada nó é medido uma única vez por quadro, o que mantém o layout rápido mesmo com milhares de nós. No Compose, cada \cod{Modifier} ajusta as restrições antes de repassá-las; no Flutter, \cod{Padding} e \cod{SizedBox} fazem esse papel.
      \end{pordentro}
    \end{column}
  \end{columns}
\end{frame}

\begin{frame}{Do estado aos pixels: o caminho de um quadro}
  \centering
  \begin{tikzpicture}[node distance=5mm,every node/.style={font=\scriptsize}]
    \node[caixaambar,minimum width=2cm,minimum height=11mm] (s) {estado\\ muda};
    \node[caixapetroleo,right=of s,minimum width=2.3cm,minimum height=11mm] (c) {\textbf{composição}\\ \tiny o que mostrar};
    \node[caixapetroleo,right=of c,minimum width=2.3cm,minimum height=11mm] (l) {\textbf{layout}\\ \tiny tamanho e posição};
    \node[caixapetroleo,right=of l,minimum width=2.3cm,minimum height=11mm] (d) {\textbf{desenho}\\ \tiny cores, textos, formas};
    \node[caixaescura,right=of d,minimum width=2.3cm,minimum height=11mm] (r) {\textbf{rasterização}\\ \tiny GPU gera os pixels};
    \draw[seta] (s) -- (c); \draw[seta] (c) -- (l); \draw[seta] (l) -- (d); \draw[seta] (d) -- (r);
    \node[rotulo,below=1mm of c] {Compose: \emph{Composition}\\ Flutter: \emph{build}};
    \node[rotulo,below=1mm of l] {\emph{Layout}\\ \emph{layout}};
    \node[rotulo,below=1mm of d] {\emph{Drawing}\\ \emph{paint} e \emph{composite}};
    \node[rotulo,below=1mm of r] {motor gráfico\\ (Skia, Impeller)};
  \end{tikzpicture}
  \raggedright
  \begin{columns}[T]
    \begin{column}{0.49\textwidth}
      \begin{definicao}[Quadro e orçamento de tempo]
        A tela é redesenhada no ritmo do sinal vertical do display (VSync): a 60\,Hz cada quadro tem \textbf{16,7\,ms}; a 120\,Hz, só \textbf{8,3\,ms}. Se o trabalho não couber, o quadro se perde e a animação engasga (\emph{jank}).
      \end{definicao}
    \end{column}
    \begin{column}{0.49\textwidth}
      \begin{pordentro}[Pular fases economiza tempo]
        Uma cor que anima só precisa de desenho; uma posição que muda, de layout e desenho. Ler o estado na fase mais tardia possível evita repetir as anteriores.
      \end{pordentro}
      \begin{dica}
        O \emph{Performance} do Flutter DevTools e o \emph{Layout Inspector} do Android Studio mostram onde o tempo do quadro foi gasto.
      \end{dica}
    \end{column}
  \end{columns}
\end{frame}

\begin{frame}{Por dentro: descrição, instância e renderização}
  \begin{columns}[T]
    \begin{column}{0.56\textwidth}
      \scriptsize
      \begin{tabularx}{\linewidth}{@{}L l l@{}}
        \toprule
        \cabecalho{Papel} & \cabecalho{Flutter} & \cabecalho{Compose} \\
        \midrule
        Descrição leve, recriada a cada mudança & \cod{Widget} & chamada \cod{@Composable} \\
        Instância duradoura, guarda o estado & \cod{Element} + \cod{State} & composição (\cod{remember}) \\
        Mede, posiciona e desenha & \cod{RenderObject} & \cod{LayoutNode} \\
        Informa leitores de tela & árvore de semântica & árvore de semântica \\
        \bottomrule
      \end{tabularx}
      \normalsize
      \begin{analogia}
        O widget é a planta do arquiteto: barata de redesenhar. O elemento é o apartamento construído, onde mora o estado. O objeto de renderização é a estrutura medida e pintada. Mudar a planta não obriga a demolir o prédio.
      \end{analogia}
    \end{column}
    \begin{column}{0.42\textwidth}
      \begin{conceito}[Por que separar]
        Descrições são criadas aos milhares por segundo, então precisam ser baratas. Medir e desenhar é caro, então os objetos de renderização são \textbf{reaproveitados} enquanto possível.
      \end{conceito}
      \begin{pordentro}[Semântica e acessibilidade]
        Em paralelo, o framework monta uma árvore de semântica (rótulos, papéis, ações) que o TalkBack e o VoiceOver leem. Ícone sem rótulo fica mudo para quem não enxerga (Módulo 04).
      \end{pordentro}
    \end{column}
  \end{columns}
\end{frame}

\begin{frame}{Atualizar sem refazer tudo}
  \begin{columns}[c]
    \begin{column}{0.42\textwidth}
      \centering
      \begin{tikzpicture}[level distance=10mm,
          every node/.style={caixa,font=\tiny,minimum height=5mm,minimum width=9mm,inner sep=2pt},
          edge from parent/.style={draw=azulpetroleo,line width=0.6pt},
          level 1/.style={sibling distance=22mm},
          level 2/.style={sibling distance=11mm}]
        \node {Tela}
          child {node {Barra} child {node {Título}}}
          child {node[caixaambarforte] (m) {Medidor}
            child {node[caixaambar] {Valor}}
            child {node[caixaambar] {Barra}}};
        \node[rotulo,below=11mm of m,text width=3.4cm] {o estado \cod{temperatura} mudou: só o \textbf{Medidor} e seus filhos são reconstruídos};
      \end{tikzpicture}
    \end{column}
    \begin{column}{0.56\textwidth}
      \begin{conceito}[Flutter: \emph{rebuild} e reconciliação]
        \cod{setState} marca o elemento como sujo; no próximo quadro, só ele executa \cod{build}. Filhos de mesmo tipo e chave são reaproveitados e só atualizados.
      \end{conceito}
      \begin{conceito}[Compose: recomposição inteligente]
        O Compose sabe quais funções leram cada estado e só executa essas de novo; as que recebem parâmetros iguais aos anteriores são puladas (\emph{skipping}).
      \end{conceito}
      \begin{atencao}[Consequência importante]
        Funções de interface podem rodar muitas vezes, em qualquer ordem: devem ser \textbf{rápidas e sem efeitos colaterais}. Rede e banco ficam em APIs de efeito (Módulo 07).
      \end{atencao}
    \end{column}
  \end{columns}
\end{frame}

\begin{frame}[fragile]{Identidade e chaves}
  \begin{columns}[T]
    \begin{column}{0.47\textwidth}
      \centering
      \begin{tikzpicture}[baseline=(current bounding box.north),every node/.style={font=\tiny}]
        \tikzset{item/.style={caixa,minimum width=1.5cm,minimum height=5mm,inner sep=1pt}}
        \node[rotulo] at (0.75,0.55) {antes};
        \node[item,caixaambar] (a) at (0.75,0) {A \quad expandido};
        \node[item] at (0.75,-0.6) {B};
        \node[item] at (0.75,-1.2) {C};
        \node[rotulo,text width=2cm] at (2.95,0.55) {remove A, \textbf{sem} chave};
        \node[item,caixaambar,draw=vermelhoerro] at (2.95,0) {B \quad expandido};
        \node[item] at (2.95,-0.6) {C};
        \node[rotulo,text width=2cm] at (5.15,0.55) {remove A, \textbf{com} chave};
        \node[item] at (5.15,0) {B};
        \node[item] at (5.15,-0.6) {C};
        \node[rotulo,text width=2cm,text=vermelhoerro] at (2.95,-1.25) {o estado ficou na posição 1};
        \node[rotulo,text width=2cm,text=verdeok] at (5.15,-1.25) {o estado saiu junto com A};
      \end{tikzpicture}
\begin{codigo}[title=Compose]{Kotlin}
items(itens, key = { it.id }) { eq ->
  ItemEquipamento(eq)
}
\end{codigo}
\begin{codigo}[title=Flutter]{Dart}
ItemEquipamento(
    key: ValueKey(eq.id), eq: eq, onTap: abrir)
\end{codigo}
    \end{column}
    \begin{column}{0.51\textwidth}
      \begin{definicao}[Chave (\emph{key})]
        Identificador estável que diz ao framework qual item é qual entre duas versões da árvore. Sem chave, a correspondência é feita pela \textbf{posição}.
      \end{definicao}
      \begin{conceito}[Quando importa]
        Em listas cujos itens têm estado próprio (expandido, campo digitado, animação) e podem ser inseridos, removidos ou reordenados.
      \end{conceito}
      \begin{dica}
        Use um identificador do domínio, como o \cod{id} do equipamento. Nunca o índice da lista nem um valor aleatório criado a cada reconstrução.
      \end{dica}
    \end{column}
  \end{columns}
\end{frame}

\begin{frame}{Temas, \emph{tokens} e contexto implícito}
  \begin{columns}[c]
    \begin{column}{0.42\textwidth}
      \centering
      \begin{tikzpicture}[node distance=7mm and 3mm,every node/.style={font=\tiny}]
        \node[caixaambarforte,minimum width=2.4cm,font=\tiny\bfseries] (t) {Tema};
        \node[caixa,below=of t,minimum width=1.6cm,font=\tiny] (s) {Tela};
        \node[caixa,below left=of s,minimum width=1.3cm,font=\tiny] (l) {Lista};
        \node[caixa,below right=of s,minimum width=1.3cm,font=\tiny] (b) {Botão};
        \node[caixa,below=of l,minimum width=1.3cm,font=\tiny] (c) {Cartão};
        \draw[seta] (t) -- (s); \draw[seta] (s) -- (l); \draw[seta] (s) -- (b); \draw[seta] (l) -- (c);
        \draw[setaambar,dashed] (t.east) -| (b.north);
        \draw[setaambar,dashed] (t.west) -| ([xshift=-3mm]c.west) -- (c.west);
        \node[rotulo,text width=4cm,below=2mm of c.south east] {tracejado: componentes leem o tema direto do contexto, sem receber parâmetros};
      \end{tikzpicture}
    \end{column}
    \begin{column}{0.56\textwidth}
      \begin{definicao}[\emph{Tokens} de design]
        Decisões visuais nomeadas pelo \textbf{papel}, não pelo valor: \cod{primary}, \cod{surface}, \cod{titleMedium}. O componente pede ``a cor primária''; o tema decide qual é.
      \end{definicao}
      \begin{pordentro}[Contexto implícito]
        Tema, idioma e tamanho da tela ficam disponíveis para toda uma subárvore: \cod{InheritedWidget} no Flutter (\cod{Theme.of(context)}), \cod{CompositionLocal} no Compose. Quem lê é atualizado quando o valor muda.
      \end{pordentro}
      \begin{dica}
        Com \emph{tokens}, tema escuro e cor dinâmica (Android 12+) funcionam sem mexer nos componentes.
      \end{dica}
    \end{column}
  \end{columns}
\end{frame}

\begin{frame}{Pixels, densidade e telas de todos os tamanhos}
  \begin{columns}[T]
    \begin{column}{0.49\textwidth}
      \centering
      \begin{tikzpicture}[baseline=(current bounding box.north)]
        \draw[cinzamedio,step=0.2] (0,0) grid (1.2,1.2);
        \fill[ambar,opacity=0.7] (0.2,0.2) rectangle (1.0,1.0);
        \node[rotulo] at (0.6,-0.25) {160 dpi: 48\,dp = 48\,px};
        \draw[cinzamedio,step=0.0667] (2.6,0) grid (3.8,1.2);
        \fill[ambar,opacity=0.7] (2.8,0.2) rectangle (3.6,1.0);
        \node[rotulo] at (3.2,-0.25) {480 dpi: 48\,dp = 144\,px};
      \end{tikzpicture}
      \raggedright
      \begin{definicao}[Pixel independente de densidade]
        \cod{px = dp} $\times$ \cod{(dpi} $\div$ \cod{160)}: um botão de 48\,dp tem o mesmo tamanho físico em qualquer tela. No Flutter, é o pixel lógico (fator em \cod{devicePixelRatio}).
      \end{definicao}
      \begin{dica}[Texto]
        Fontes também seguem a escala escolhida pelo usuário: \cod{sp} no Android, \cod{TextScaler} no Flutter.
      \end{dica}
    \end{column}
    \begin{column}{0.49\textwidth}
      \begin{conceito}[Responsivo $\times$ adaptativo]
        \textbf{Responsivo}: o mesmo layout se estica ou reorganiza. \textbf{Adaptativo}: um layout por classe de tamanho de janela (barra inferior no celular, trilho lateral no tablet). Pense em \textbf{espaço disponível}, não em modelo de aparelho.
      \end{conceito}
      \begin{atencao}[Android 16 e telas grandes]
        Com \cod{targetSdk} 36 (exigido pela Play desde agosto de 2026), telas com largura mínima $\geq$ 600\,dp ignoram travas de orientação e de redimensionamento; no API 37 não haverá exceção.
      \end{atencao}
    \end{column}
  \end{columns}
\end{frame}

% =============================================================================
\section{Fundamentos: interfaces declarativas}

\begin{frame}{Interfaces imperativas e declarativas}
  \begin{columns}[T]
    \begin{column}{0.49\textwidth}
      \begin{definicao}[Interface imperativa]
        O programa cria os componentes uma vez e depois os \textbf{altera} passo a passo (\cod{setText}, \cod{setVisibility}). Basta esquecer um deles para a tela mostrar dados errados.
      \end{definicao}
    \end{column}
    \begin{column}{0.49\textwidth}
      \begin{definicao}[Interface declarativa]
        O programa \textbf{descreve} como a tela deve ser para cada estado. Quando o estado muda, o framework gera a nova descrição e aplica apenas as diferenças.
      \end{definicao}
    \end{column}
  \end{columns}
  \vspace{2pt}
  \centering
  \begin{tikzpicture}[node distance=2mm and 6mm,every node/.style={font=\scriptsize}]
    \node[rotulo,anchor=east,text width=1.6cm,align=right] (li) at (0,0) {\textbf{imperativo}};
    \node[caixa,right=3mm of li,minimum width=1.6cm] (e1) {estado muda};
    \node[caixaambar,right=of e1,minimum width=3.2cm] (c1) {seu código acha e altera\\ cada componente afetado};
    \node[caixa,right=of c1,minimum width=2.2cm] (t1) {tela (mutável)};
    \draw[seta] (e1) -- (c1); \draw[seta] (c1) -- (t1);
    \node[rotulo,anchor=east,text width=1.6cm,align=right] (ld) at (0,-1.05) {\textbf{declarativo}};
    \node[caixa,right=3mm of ld,minimum width=1.6cm] (e2) {estado muda};
    \node[caixapetroleo,right=of e2,minimum width=2.1cm] (f2) {UI = f(estado)};
    \node[caixa,right=of f2,minimum width=2.1cm] (d2) {nova descrição};
    \node[caixapetroleo,right=of d2,minimum width=2.6cm] (a2) {framework compara\\ e aplica diferenças};
    \draw[seta] (e2) -- (f2); \draw[seta] (f2) -- (d2); \draw[seta] (d2) -- (a2);
  \end{tikzpicture}
  \raggedright
  \begin{historia}
    Views do Android e UIKit do iOS (2008) são imperativos. O React (2013) popularizou a UI declarativa, depois adotada pelo Flutter (2018), SwiftUI (2019) e Jetpack Compose (2021).
  \end{historia}
\end{frame}

\begin{frame}{Composição: a tela como uma árvore}
  \begin{columns}[c]
    \begin{column}{0.52\textwidth}
      \centering
      \begin{tikzpicture}[level distance=10mm,
          every node/.style={caixa,font=\tiny,minimum height=5mm,inner sep=2pt},
          edge from parent/.style={draw=azulpetroleo,line width=0.6pt},
          level 1/.style={sibling distance=30mm},
          level 2/.style={sibling distance=21mm},
          level 3/.style={sibling distance=10mm}]
        \node[caixaescura,font=\tiny] (raiz) {TelaEquipamentos}
          child {node {BarraSuperior}}
          child {node {ListaEquipamentos}
            child {node[caixaambar,font=\tiny] {CartaoBancada}
              child {node {Ícone}}
              child {node {Textos}}
              child {node {Barra}}}
            child {node[caixaambar,font=\tiny] {CartaoBancada}}};
        \draw[seta] (-2.6,0.2) -- (-2.6,-3.2);
        \draw[setaambar] (-2.2,-3.2) -- (-2.2,0.2);
        \node[rotulo,anchor=north west,text width=6cm,align=left] at (-2.75,-3.45) {\textcolor{azulpetroleo}{$\downarrow$} dados descem por parâmetros\quad \textcolor{ambar}{$\uparrow$} eventos sobem por \emph{callbacks}};
      \end{tikzpicture}
    \end{column}
    \begin{column}{0.46\textwidth}
      \begin{definicao}[Composição]
        Construir interfaces complexas \textbf{combinando} componentes pequenos e independentes, em vez de herdar de componentes grandes. Cada componente recebe dados e devolve uma descrição de UI.
      \end{definicao}
      \begin{analogia}
        Como peças de montar: o cartão é feito de ícone, textos e indicador; a lista, de vários cartões. Trocar uma peça não obriga a refazer as outras.
      \end{analogia}
      \begin{dica}[Um bom componente]
        Faz uma coisa só, recebe tudo por parâmetro, não sabe de onde vêm os dados e pode ser visto isoladamente em uma \emph{preview}.
      \end{dica}
    \end{column}
  \end{columns}
\end{frame}

\begin{frame}{Layout por restrições}
  \begin{columns}[c]
    \begin{column}{0.47\textwidth}
      \centering
      \begin{tikzpicture}[node distance=11mm]
        \node[caixaescura,minimum width=3.4cm] (t) {Tela \quad 360 $\times$ 800};
        \node[caixa,minimum width=3.4cm,below=of t] (c) {Coluna};
        \node[caixaambar,minimum width=3.4cm,below=of c] (k) {Cartão};
        \draw[seta] ([xshift=-10mm]t.south) -- node[rotulo,left,text width=2.1cm,align=right]{largura 0--360\\ altura 0--800} ([xshift=-10mm]c.north);
        \draw[seta] ([xshift=-10mm]c.south) -- node[rotulo,left,text width=2.1cm,align=right]{largura 0--360\\ altura livre ($\infty$)} ([xshift=-10mm]k.north);
        \draw[setaambar] ([xshift=10mm]k.north) -- node[rotulo,right,text width=1.6cm,align=left]{``tenho\\ 360 $\times$ 96''} ([xshift=10mm]c.south);
        \node[rotulo,below=1.5mm of k,text width=4cm] {a coluna põe o cartão em $y = 0$\\ e o próximo em $y = 96$};
      \end{tikzpicture}
    \end{column}
    \begin{column}{0.51\textwidth}
      \begin{definicao}[Restrições (\emph{constraints})]
        Intervalo de tamanhos que o pai aceita para o filho: largura e altura mínimas e máximas. \textbf{Justa} quando mínimo = máximo; \textbf{frouxa} quando o mínimo é zero.
      \end{definicao}
      \begin{conceito}[O algoritmo em três frases]
        \textbf{Restrições descem}: o pai diz os limites. \textbf{Tamanhos sobem}: o filho escolhe um tamanho dentro deles. \textbf{O pai posiciona}: o filho não sabe onde está.
      \end{conceito}
      \begin{pordentro}[Uma passada só]
        Cada nó é medido uma única vez por quadro, o que mantém o layout rápido mesmo com milhares de nós. No Compose, cada \cod{Modifier} ajusta as restrições antes de repassá-las; no Flutter, \cod{Padding} e \cod{SizedBox} fazem esse papel.
      \end{pordentro}
    \end{column}
  \end{columns}
\end{frame}

\begin{frame}{Do estado aos pixels: o caminho de um quadro}
  \centering
  \begin{tikzpicture}[node distance=5mm,every node/.style={font=\scriptsize}]
    \node[caixaambar,minimum width=2cm,minimum height=11mm] (s) {estado\\ muda};
    \node[caixapetroleo,right=of s,minimum width=2.3cm,minimum height=11mm] (c) {\textbf{composição}\\ \tiny o que mostrar};
    \node[caixapetroleo,right=of c,minimum width=2.3cm,minimum height=11mm] (l) {\textbf{layout}\\ \tiny tamanho e posição};
    \node[caixapetroleo,right=of l,minimum width=2.3cm,minimum height=11mm] (d) {\textbf{desenho}\\ \tiny cores, textos, formas};
    \node[caixaescura,right=of d,minimum width=2.3cm,minimum height=11mm] (r) {\textbf{rasterização}\\ \tiny GPU gera os pixels};
    \draw[seta] (s) -- (c); \draw[seta] (c) -- (l); \draw[seta] (l) -- (d); \draw[seta] (d) -- (r);
    \node[rotulo,below=1mm of c] {Compose: \emph{Composition}\\ Flutter: \emph{build}};
    \node[rotulo,below=1mm of l] {\emph{Layout}\\ \emph{layout}};
    \node[rotulo,below=1mm of d] {\emph{Drawing}\\ \emph{paint} e \emph{composite}};
    \node[rotulo,below=1mm of r] {motor gráfico\\ (Skia, Impeller)};
  \end{tikzpicture}
  \raggedright
  \begin{columns}[T]
    \begin{column}{0.49\textwidth}
      \begin{definicao}[Quadro e orçamento de tempo]
        A tela é redesenhada no ritmo do sinal vertical do display (VSync): a 60\,Hz cada quadro tem \textbf{16,7\,ms}; a 120\,Hz, só \textbf{8,3\,ms}. Se o trabalho não couber, o quadro se perde e a animação engasga (\emph{jank}).
      \end{definicao}
    \end{column}
    \begin{column}{0.49\textwidth}
      \begin{pordentro}[Pular fases economiza tempo]
        Uma cor que anima só precisa de desenho; uma posição que muda, de layout e desenho. Ler o estado na fase mais tardia possível evita repetir as anteriores.
      \end{pordentro}
      \begin{dica}
        O \emph{Performance} do Flutter DevTools e o \emph{Layout Inspector} do Android Studio mostram onde o tempo do quadro foi gasto.
      \end{dica}
    \end{column}
  \end{columns}
\end{frame}

\begin{frame}{Por dentro: descrição, instância e renderização}
  \begin{columns}[T]
    \begin{column}{0.56\textwidth}
      \scriptsize
      \begin{tabularx}{\linewidth}{@{}L l l@{}}
        \toprule
        \cabecalho{Papel} & \cabecalho{Flutter} & \cabecalho{Compose} \\
        \midrule
        Descrição leve, recriada a cada mudança & \cod{Widget} & chamada \cod{@Composable} \\
        Instância duradoura, guarda o estado & \cod{Element} + \cod{State} & composição (\cod{remember}) \\
        Mede, posiciona e desenha & \cod{RenderObject} & \cod{LayoutNode} \\
        Informa leitores de tela & árvore de semântica & árvore de semântica \\
        \bottomrule
      \end{tabularx}
      \normalsize
      \begin{analogia}
        O widget é a planta do arquiteto: barata de redesenhar. O elemento é o apartamento construído, onde mora o estado. O objeto de renderização é a estrutura medida e pintada. Mudar a planta não obriga a demolir o prédio.
      \end{analogia}
    \end{column}
    \begin{column}{0.42\textwidth}
      \begin{conceito}[Por que separar]
        Descrições são criadas aos milhares por segundo, então precisam ser baratas. Medir e desenhar é caro, então os objetos de renderização são \textbf{reaproveitados} enquanto possível.
      \end{conceito}
      \begin{pordentro}[Semântica e acessibilidade]
        Em paralelo, o framework monta uma árvore de semântica (rótulos, papéis, ações) que o TalkBack e o VoiceOver leem. Ícone sem rótulo fica mudo para quem não enxerga (Módulo 04).
      \end{pordentro}
    \end{column}
  \end{columns}
\end{frame}

\begin{frame}{Atualizar sem refazer tudo}
  \begin{columns}[c]
    \begin{column}{0.42\textwidth}
      \centering
      \begin{tikzpicture}[level distance=10mm,
          every node/.style={caixa,font=\tiny,minimum height=5mm,minimum width=9mm,inner sep=2pt},
          edge from parent/.style={draw=azulpetroleo,line width=0.6pt},
          level 1/.style={sibling distance=22mm},
          level 2/.style={sibling distance=11mm}]
        \node {Tela}
          child {node {Barra} child {node {Título}}}
          child {node[caixaambarforte] (m) {Medidor}
            child {node[caixaambar] {Valor}}
            child {node[caixaambar] {Barra}}};
        \node[rotulo,below=11mm of m,text width=3.4cm] {o estado \cod{temperatura} mudou: só o \textbf{Medidor} e seus filhos são reconstruídos};
      \end{tikzpicture}
    \end{column}
    \begin{column}{0.56\textwidth}
      \begin{conceito}[Flutter: \emph{rebuild} e reconciliação]
        \cod{setState} marca o elemento como sujo; no próximo quadro, só ele executa \cod{build}. Filhos de mesmo tipo e chave são reaproveitados e só atualizados.
      \end{conceito}
      \begin{conceito}[Compose: recomposição inteligente]
        O Compose sabe quais funções leram cada estado e só executa essas de novo; as que recebem parâmetros iguais aos anteriores são puladas (\emph{skipping}).
      \end{conceito}
      \begin{atencao}[Consequência importante]
        Funções de interface podem rodar muitas vezes, em qualquer ordem: devem ser \textbf{rápidas e sem efeitos colaterais}. Rede e banco ficam em APIs de efeito (Módulo 07).
      \end{atencao}
    \end{column}
  \end{columns}
\end{frame}

\begin{frame}[fragile]{Identidade e chaves}
  \begin{columns}[T]
    \begin{column}{0.47\textwidth}
      \centering
      \begin{tikzpicture}[baseline=(current bounding box.north),every node/.style={font=\tiny}]
        \tikzset{item/.style={caixa,minimum width=1.5cm,minimum height=5mm,inner sep=1pt}}
        \node[rotulo] at (0.75,0.55) {antes};
        \node[item,caixaambar] (a) at (0.75,0) {A \quad expandido};
        \node[item] at (0.75,-0.6) {B};
        \node[item] at (0.75,-1.2) {C};
        \node[rotulo,text width=2cm] at (2.95,0.55) {remove A, \textbf{sem} chave};
        \node[item,caixaambar,draw=vermelhoerro] at (2.95,0) {B \quad expandido};
        \node[item] at (2.95,-0.6) {C};
        \node[rotulo,text width=2cm] at (5.15,0.55) {remove A, \textbf{com} chave};
        \node[item] at (5.15,0) {B};
        \node[item] at (5.15,-0.6) {C};
        \node[rotulo,text width=2cm,text=vermelhoerro] at (2.95,-1.25) {o estado ficou na posição 1};
        \node[rotulo,text width=2cm,text=verdeok] at (5.15,-1.25) {o estado saiu junto com A};
      \end{tikzpicture}
\begin{codigo}[title=Compose]{Kotlin}
items(itens, key = { it.id }) { eq ->
  ItemEquipamento(eq)
}
\end{codigo}
\begin{codigo}[title=Flutter]{Dart}
ItemEquipamento(
    key: ValueKey(eq.id), eq: eq, onTap: abrir)
\end{codigo}
    \end{column}
    \begin{column}{0.51\textwidth}
      \begin{definicao}[Chave (\emph{key})]
        Identificador estável que diz ao framework qual item é qual entre duas versões da árvore. Sem chave, a correspondência é feita pela \textbf{posição}.
      \end{definicao}
      \begin{conceito}[Quando importa]
        Em listas cujos itens têm estado próprio (expandido, campo digitado, animação) e podem ser inseridos, removidos ou reordenados.
      \end{conceito}
      \begin{dica}
        Use um identificador do domínio, como o \cod{id} do equipamento. Nunca o índice da lista nem um valor aleatório criado a cada reconstrução.
      \end{dica}
    \end{column}
  \end{columns}
\end{frame}

\begin{frame}{Temas, \emph{tokens} e contexto implícito}
  \begin{columns}[c]
    \begin{column}{0.42\textwidth}
      \centering
      \begin{tikzpicture}[node distance=7mm and 3mm,every node/.style={font=\tiny}]
        \node[caixaambarforte,minimum width=2.4cm,font=\tiny\bfseries] (t) {Tema};
        \node[caixa,below=of t,minimum width=1.6cm,font=\tiny] (s) {Tela};
        \node[caixa,below left=of s,minimum width=1.3cm,font=\tiny] (l) {Lista};
        \node[caixa,below right=of s,minimum width=1.3cm,font=\tiny] (b) {Botão};
        \node[caixa,below=of l,minimum width=1.3cm,font=\tiny] (c) {Cartão};
        \draw[seta] (t) -- (s); \draw[seta] (s) -- (l); \draw[seta] (s) -- (b); \draw[seta] (l) -- (c);
        \draw[setaambar,dashed] (t.east) -| (b.north);
        \draw[setaambar,dashed] (t.west) -| ([xshift=-3mm]c.west) -- (c.west);
        \node[rotulo,text width=4cm,below=2mm of c.south east] {tracejado: componentes leem o tema direto do contexto, sem receber parâmetros};
      \end{tikzpicture}
    \end{column}
    \begin{column}{0.56\textwidth}
      \begin{definicao}[\emph{Tokens} de design]
        Decisões visuais nomeadas pelo \textbf{papel}, não pelo valor: \cod{primary}, \cod{surface}, \cod{titleMedium}. O componente pede ``a cor primária''; o tema decide qual é.
      \end{definicao}
      \begin{pordentro}[Contexto implícito]
        Tema, idioma e tamanho da tela ficam disponíveis para toda uma subárvore: \cod{InheritedWidget} no Flutter (\cod{Theme.of(context)}), \cod{CompositionLocal} no Compose. Quem lê é atualizado quando o valor muda.
      \end{pordentro}
      \begin{dica}
        Com \emph{tokens}, tema escuro e cor dinâmica (Android 12+) funcionam sem mexer nos componentes.
      \end{dica}
    \end{column}
  \end{columns}
\end{frame}

\begin{frame}{Pixels, densidade e telas de todos os tamanhos}
  \begin{columns}[T]
    \begin{column}{0.49\textwidth}
      \centering
      \begin{tikzpicture}[baseline=(current bounding box.north)]
        \draw[cinzamedio,step=0.2] (0,0) grid (1.2,1.2);
        \fill[ambar,opacity=0.7] (0.2,0.2) rectangle (1.0,1.0);
        \node[rotulo] at (0.6,-0.25) {160 dpi: 48\,dp = 48\,px};
        \draw[cinzamedio,step=0.0667] (2.6,0) grid (3.8,1.2);
        \fill[ambar,opacity=0.7] (2.8,0.2) rectangle (3.6,1.0);
        \node[rotulo] at (3.2,-0.25) {480 dpi: 48\,dp = 144\,px};
      \end{tikzpicture}
      \raggedright
      \begin{definicao}[Pixel independente de densidade]
        \cod{px = dp} $\times$ \cod{(dpi} $\div$ \cod{160)}: um botão de 48\,dp tem o mesmo tamanho físico em qualquer tela. No Flutter, é o pixel lógico (fator em \cod{devicePixelRatio}).
      \end{definicao}
      \begin{dica}[Texto]
        Fontes também seguem a escala escolhida pelo usuário: \cod{sp} no Android, \cod{TextScaler} no Flutter.
      \end{dica}
    \end{column}
    \begin{column}{0.49\textwidth}
      \begin{conceito}[Responsivo $\times$ adaptativo]
        \textbf{Responsivo}: o mesmo layout se estica ou reorganiza. \textbf{Adaptativo}: um layout por classe de tamanho de janela (barra inferior no celular, trilho lateral no tablet). Pense em \textbf{espaço disponível}, não em modelo de aparelho.
      \end{conceito}
      \begin{atencao}[Android 16 e telas grandes]
        Com \cod{targetSdk} 36 (exigido pela Play desde agosto de 2026), telas com largura mínima $\geq$ 600\,dp ignoram travas de orientação e de redimensionamento; no API 37 não haverá exceção.
      \end{atencao}
    \end{column}
  \end{columns}
\end{frame}


% =============================================================================
\section{Declarativo na prática}

\begin{frame}{Do imperativo ao Compose}
  \begin{columns}[T]
    \begin{column}{0.48\textwidth}
      \begin{exemplo}[Imperativo (Views + XML)]
        \small
        Você cria a árvore de \emph{widgets} e a \textbf{modifica} manualmente quando os dados mudam:\\[2pt]
        \cod{texto.setText("25 °C")}\\
        \cod{alerta.setVisibility(View.VISIBLE)}\\[2pt]
        Fácil esquecer algo e deixar a tela inconsistente.
      \end{exemplo}
    \end{column}
    \begin{column}{0.48\textwidth}
      \begin{exemplo}[Declarativo (Compose)]
        \small
        Você descreve \textbf{como a tela deve ser} para um dado estado. Quando o estado muda, o Compose \textbf{redesenha} o que for necessário:\\[2pt]
        \cod{Text("\$temp °C")}\\
        \cod{if (temp > 30) Alerta()}\\[2pt]
        A tela fica sempre consistente.
      \end{exemplo}
    \end{column}
  \end{columns}
  \begin{center}
  \begin{tikzpicture}[node distance=10mm]
    \node[caixaambar,minimum width=2.6cm] (e) {\textbf{Estado}\\ \cod{temp = 31.2}};
    \node[caixaescura,minimum width=2.6cm,right=of e] (f) {\textbf{Funções}\\ \cod{@Composable}};
    \node[caixa,minimum width=2.6cm,right=of f] (u) {\textbf{UI}\\ na tela};
    \draw[seta] (e) -- node[rotulo,above]{entrada} (f);
    \draw[seta] (f) -- node[rotulo,above]{emite} (u);
    \draw[setaambar,rounded corners] (u.south) -- ++(0,-5mm) -| node[rotulo,pos=0.25,below]{evento do usuário altera o estado} (e.south);
  \end{tikzpicture}
  \end{center}
  \begin{destaque}
    \textbf{UI = f(estado)}
  \end{destaque}
\end{frame}

\begin{frame}[fragile]{Funções \cod{@Composable}}
  \begin{columns}[T]
    \begin{column}{0.52\textwidth}
\begin{codigo}[title=CartaoBancada.kt]{Kotlin}
@Composable
fun CartaoBancada(
  numero: Int,
  temperatura: Double,
  modifier: Modifier = Modifier
) {
  Column(modifier = modifier.padding(16.dp)) {
    Text(
      text = "Bancada $numero",
      style = MaterialTheme.typography.titleMedium
    )
    Text("%.1f °C".format(temperatura))
    if (temperatura > 30.0) {
      Text("Acima do limite!", color = Color.Red)
    }
  }
}
\end{codigo}
    \end{column}
    \begin{column}{0.46\textwidth}
      \small
      \begin{itemize}
        \item São funções Kotlin comuns anotadas com \cod{@Composable}.
        \item Nome em \textbf{PascalCase} (são ``substantivos'').
        \item Não retornam nada: \textbf{emitem} elementos de UI.
        \item Podem usar \cod{if}, \cod{for} e \cod{when} livremente.
        \item Recebem um \cod{modifier} opcional como primeiro parâmetro com valor padrão.
      \end{itemize}
      \begin{atencao}[Devem ser rápidas e sem efeitos colaterais]
        O Compose pode chamá-las muitas vezes, em qualquer ordem. Nada de rede, banco ou \cod{Log} pesado dentro delas.
      \end{atencao}
    \end{column}
  \end{columns}
\end{frame}

% =============================================================================
\section{Componentes básicos}

\begin{frame}[fragile]{Texto, botões, ícones e imagens}
  \begin{columns}[T]
    \begin{column}{0.54\textwidth}
\begin{codigo}[title=Componentes]{Kotlin}
Text(
  text = stringResource(R.string.titulo),
  style = MaterialTheme.typography.headlineSmall,
  maxLines = 1,
  overflow = TextOverflow.Ellipsis
)

Button(onClick = { reservar() }) { Text("Reservar") }
OutlinedButton(onClick = {}) { Text("Cancelar") }
TextButton(onClick = {}) { Text("Detalhes") }
IconButton(onClick = { abrirMenu() }) {
  Icon(Icons.Default.MoreVert, contentDescription = "Menu")
}

Image(
  painter = painterResource(R.drawable.osciloscopio),
  contentDescription = "Foto do osciloscópio",
  contentScale = ContentScale.Crop,
  modifier = Modifier.size(64.dp).clip(CircleShape)
)
\end{codigo}
    \end{column}
    \begin{column}{0.44\textwidth}
      \centering
      \begin{tikzpicture}[font=\scriptsize]
        \draw[petroleo,rounded corners=3pt,fill=petroleo!4] (0,0) rectangle (4.6,3.5);
        \node[anchor=west,font=\small\bfseries,text=cinzatexto] at (0.2,3.1) {Bancada 3 -- Osciloscópio};
        \fill[petroleo,rounded corners=8pt] (0.2,2.05) rectangle (1.6,2.55) node[midway,white] {Reservar};
        \draw[petroleo,rounded corners=8pt] (1.75,2.05) rectangle (3.15,2.55) node[midway,petroleo] {Cancelar};
        \node[text=petroleo] at (3.85,2.3) {Detalhes};
        \fill[ambar!60] (0.75,1.05) circle (0.45);
        \node[font=\tiny,white] at (0.75,1.05) {img};
        \node[anchor=west,text width=2.9cm,font=\tiny,text=cinzamedio] at (1.4,1.05) {\cod{Image} recortada em círculo com \cod{clip(CircleShape)}};
        \node[text=cinzamedio] at (4.3,3.1) {$\vdots$};
        \node[rotulo] at (2.3,0.3) {resultado aproximado};
      \end{tikzpicture}
      \begin{atencao}[Acessibilidade]
        Toda imagem informativa precisa de \cod{contentDescription}; imagens decorativas usam \cod{null}.
      \end{atencao}
    \end{column}
  \end{columns}
\end{frame}

\begin{frame}[fragile]{Campos de entrada, seleção e cartões}
  \begin{columns}[T]
    \begin{column}{0.54\textwidth}
\begin{codigo}[title=Entrada de dados]{Kotlin}
var descricao by remember { mutableStateOf("") }
var ligado by remember { mutableStateOf(false) }

OutlinedTextField(
  value = descricao,
  onValueChange = { descricao = it },
  label = { Text("Descrição") },
  singleLine = true,
  isError = descricao.isBlank(),
  keyboardOptions = KeyboardOptions(
    keyboardType = KeyboardType.Text)
)

Row(verticalAlignment = Alignment.CenterVertically) {
  Text("Alimentação")
  Switch(checked = ligado, onCheckedChange = { ligado = it })
}

Card(modifier = Modifier.fillMaxWidth()) {
  Text("Conteúdo do cartão", Modifier.padding(16.dp))
}
\end{codigo}
    \end{column}
    \begin{column}{0.44\textwidth}
      \begin{conceito}[Componentes controlados]
        O \cod{TextField} \textbf{não guarda} o texto: ele mostra o \cod{value} recebido e avisa mudanças por \cod{onValueChange}. Quem guarda o estado é você (Módulo 07).
      \end{conceito}
      \scriptsize
      \begin{tabularx}{\linewidth}{@{}l L@{}}
        \toprule
        \cabecalho{Componente} & \cabecalho{Uso} \\
        \midrule
        \cod{Checkbox}, \cod{Switch} & valores booleanos \\
        \cod{RadioButton} & uma opção entre várias \\
        \cod{Slider} & valor numérico contínuo \\
        \cod{DropdownMenu} & menus e listas suspensas \\
        \cod{AlertDialog} & confirmações \\
        \cod{CircularProgressIndicator} & carregamento \\
        \bottomrule
      \end{tabularx}
    \end{column}
  \end{columns}
\end{frame}

% =============================================================================
\section{Layouts}

\begin{frame}{\cod{Column}, \cod{Row} e \cod{Box}}
  \begin{columns}[T]
    \begin{column}{0.32\textwidth}
      \centering
      \begin{tikzpicture}
        \draw[petroleo,thick,rounded corners=2pt] (0,0) rectangle (2.6,2.6);
        \foreach \y/\t in {2.0/A,1.25/B,0.5/C}{\node[caixaescura,minimum width=2cm,minimum height=5mm] at (1.3,\y) {\t};}
        \draw[setaambar] (2.85,2.3) -- (2.85,0.3);
      \end{tikzpicture}\\[3pt]
      \textbf{\cod{Column}}\\ \footnotesize empilha na vertical
    \end{column}
    \begin{column}{0.32\textwidth}
      \centering
      \begin{tikzpicture}
        \draw[petroleo,thick,rounded corners=2pt] (0,0) rectangle (2.6,2.6);
        \foreach \x/\t in {0.45/A,1.3/B,2.15/C}{\node[caixaescura,minimum width=6mm,minimum height=2cm] at (\x,1.3) {\t};}
        \draw[setaambar] (0.2,-0.25) -- (2.4,-0.25);
      \end{tikzpicture}\\[3pt]
      \textbf{\cod{Row}}\\ \footnotesize alinha na horizontal
    \end{column}
    \begin{column}{0.32\textwidth}
      \centering
      \begin{tikzpicture}
        \draw[petroleo,thick,rounded corners=2pt] (0,0) rectangle (2.6,2.6);
        \node[caixa,minimum width=2.2cm,minimum height=2.2cm,fill=petroleo!20] at (1.3,1.3) {};
        \node[caixaescura,minimum width=1.3cm,minimum height=1.1cm] at (1.3,1.3) {B};
        \node[caixaambarforte,minimum width=5mm,minimum height=5mm] at (2.15,2.15) {C};
      \end{tikzpicture}\\[3pt]
      \textbf{\cod{Box}}\\ \footnotesize sobrepõe (eixo Z)
    \end{column}
  \end{columns}
  \vspace{6pt}
  \begin{columns}[T]
    \begin{column}{0.48\textwidth}
      \begin{conceito}[Arrangement (eixo principal)]
        \cod{Top}, \cod{Center}, \cod{Bottom}, \cod{SpaceBetween}, \cod{SpaceEvenly}, \cod{spacedBy(8.dp)}.
      \end{conceito}
    \end{column}
    \begin{column}{0.48\textwidth}
      \begin{conceito}[Alignment (eixo cruzado)]
        Em \cod{Column}: \cod{Start}, \cod{CenterHorizontally}, \cod{End}. Em \cod{Row}: \cod{Top}, \cod{CenterVertically}, \cod{Bottom}.
      \end{conceito}
    \end{column}
  \end{columns}
\end{frame}

\begin{frame}[fragile]{Compondo uma tela real}
  \begin{columns}[T]
    \begin{column}{0.58\textwidth}
\begin{codigo}[title=ItemEquipamento.kt]{Kotlin}
@Composable
fun ItemEquipamento(eq: Equipamento, onClick: () -> Unit) {
  Card(onClick = onClick, modifier = Modifier.fillMaxWidth()) {
    Row(
      modifier = Modifier.padding(12.dp),
      verticalAlignment = Alignment.CenterVertically,
      horizontalArrangement = Arrangement.spacedBy(12.dp)
    ) {
      Icon(Icons.Default.Memory, contentDescription = null)
      Column(Modifier.weight(1f)) {   // ocupa o que sobrar
        Text(eq.descricao, fontWeight = FontWeight.Bold)
        Text(eq.patrimonio,
          style = MaterialTheme.typography.bodySmall)
      }
      Box(
        Modifier.size(12.dp).clip(CircleShape)
          .background(if (eq.disponivel) Verde else Ambar)
      )
    }
  }
}
\end{codigo}
    \end{column}
    \begin{column}{0.40\textwidth}
      \centering
      \begin{tikzpicture}[font=\scriptsize]
        \draw[petroleo,rounded corners=4pt,fill=white] (0,0) rectangle (4.4,1.3);
        \draw[petroleo!40,dashed] (0.15,0.15) rectangle (4.25,1.15);
        \node[caixaescura,minimum size=6mm,inner sep=0pt] at (0.55,0.65) {\tiny CPU};
        \node[anchor=west,font=\scriptsize\bfseries] at (1.0,0.85) {Osciloscópio};
        \node[anchor=west,font=\tiny,text=cinzamedio] at (1.0,0.45) {IFMA-0042};
        \fill[verdeok] (3.95,0.65) circle (0.1);
        \draw[ambar,decorate,decoration={brace,amplitude=3pt,mirror}] (1.0,-0.05) -- (3.6,-0.05) node[midway,below=3pt,rotulo]{\cod{weight(1f)}};
        \node[rotulo] at (2.2,1.6) {\cod{Card} $\rightarrow$ \cod{Row} $\rightarrow$ [\cod{Icon} | \cod{Column} | \cod{Box}]};
      \end{tikzpicture}
      \begin{dica}
        \cod{Modifier.weight(1f)} distribui o espaço restante dentro de \cod{Row}/\cod{Column}, como o \emph{flex} do CSS.
      \end{dica}
      \begin{dica}
        \cod{Spacer(Modifier.height(8.dp))} cria espaço vazio entre elementos.
      \end{dica}
    \end{column}
  \end{columns}
\end{frame}

% =============================================================================
\section{Modifiers}

\begin{frame}[fragile]{Modifiers: a ordem importa}
  \begin{columns}[T]
    \begin{column}{0.5\textwidth}
\begin{codigo}[title={Mesmos modifiers, ordens diferentes}]{Kotlin}
// (a) fundo e DEPOIS espaçamento interno
Text("A", Modifier
  .background(Ambar)
  .padding(16.dp))

// (b) espaçamento e DEPOIS fundo
Text("B", Modifier
  .padding(16.dp)
  .background(Ambar))

// (c) clique antes do padding:
// a área de toque inclui o padding
Text("C", Modifier
  .clickable { /* ... */ }
  .padding(16.dp))
\end{codigo}
    \end{column}
    \begin{column}{0.48\textwidth}
      \centering
      \begin{tikzpicture}
        \fill[ambar] (0,0) rectangle (1.6,1.0); \node at (0.8,0.5) {\textbf{A}};
        \node[rotulo] at (0.8,-0.3) {(a) fundo inclui o padding};
        \draw[cinzamedio,dashed] (2.4,0) rectangle (4.0,1.0);
        \fill[ambar] (3.0,0.35) rectangle (3.4,0.65); \node at (3.2,0.5) {\textbf{B}};
        \node[rotulo] at (3.2,-0.3) {(b) fundo só no conteúdo};
      \end{tikzpicture}
      \begin{conceito}[Como pensar]
        Os modifiers são aplicados \textbf{em cadeia, de fora para dentro}: cada um envolve o resultado dos seguintes.
      \end{conceito}
      \scriptsize
      \renewcommand{\arraystretch}{1.1}
      \begin{tabular}{@{}l l@{}}
        \toprule
        \cabecalho{Categoria} & \cabecalho{Exemplos} \\
        \midrule
        Tamanho & \cod{size}, \cod{fillMaxWidth}, \cod{weight} \\
        Espaço & \cod{padding}, \cod{offset} \\
        Desenho & \cod{background}, \cod{border}, \cod{clip} \\
        Interação & \cod{clickable}, \cod{toggleable} \\
        Semântica & \cod{semantics}, \cod{testTag} \\
        \bottomrule
      \end{tabular}
    \end{column}
  \end{columns}
\end{frame}

% =============================================================================
\section{Listas}

\begin{frame}[fragile]{\cod{LazyColumn}: listas eficientes}
  \begin{columns}[T]
    \begin{column}{0.54\textwidth}
\begin{codigo}[title=ListaEquipamentos.kt]{Kotlin}
@Composable
fun ListaEquipamentos(
  itens: List<Equipamento>,
  onItemClick: (Equipamento) -> Unit,
  modifier: Modifier = Modifier
) {
  LazyColumn(
    modifier = modifier,
    contentPadding = PaddingValues(16.dp),
    verticalArrangement = Arrangement.spacedBy(8.dp)
  ) {
    item { Text("Equipamentos (${itens.size})") }
    items(
      items = itens,
      key = { it.patrimonio }      // chave estável!
    ) { eq ->
      ItemEquipamento(eq, onClick = { onItemClick(eq) })
    }
  }
}
\end{codigo}
    \end{column}
    \begin{column}{0.44\textwidth}
      \begin{conceito}[Lazy = sob demanda]
        Só compõe os itens visíveis na tela (mais uma margem), como a antiga \cod{RecyclerView}. Uma \cod{Column} com 1000 itens compõe todos de uma vez.
      \end{conceito}
      \begin{dica}[Use \cod{key}]
        Uma chave estável preserva o estado e as animações de cada item quando a lista é reordenada.
      \end{dica}
      \scriptsize
      \begin{tabularx}{\linewidth}{@{}l L@{}}
        \toprule
        \cod{LazyRow} & lista horizontal \\
        \cod{LazyVerticalGrid} & grade (\cod{GridCells.Adaptive(160.dp)}) \\
        \cod{stickyHeader} & cabeçalhos fixos de seção \\
        \bottomrule
      \end{tabularx}
    \end{column}
  \end{columns}
\end{frame}

% =============================================================================
\section{Material 3 e temas}

\begin{frame}[fragile]{\cod{Scaffold}: o esqueleto da tela}
  \begin{columns}[T]
    \begin{column}{0.56\textwidth}
\begin{codigo}[title=TelaEquipamentos.kt]{Kotlin}
@OptIn(ExperimentalMaterial3Api::class)
@Composable
fun TelaEquipamentos(itens: List<Equipamento>) {
  Scaffold(
    topBar = {
      TopAppBar(title = { Text("LabMonitor") })
    },
    floatingActionButton = {
      FloatingActionButton(onClick = { /* novo */ }) {
        Icon(Icons.Default.Add, "Adicionar")
      }
    },
    bottomBar = { BarraNavegacao() }
  ) { innerPadding ->
    ListaEquipamentos(
      itens,
      onItemClick = {},
      modifier = Modifier.padding(innerPadding)
    )
  }
}
\end{codigo}
    \end{column}
    \begin{column}{0.42\textwidth}
      \centering
      \begin{tikzpicture}[font=\tiny]
        \draw[petroleo,thick,rounded corners=5pt] (0,0) rectangle (2.6,4.4);
        \fill[petroleo] (0.08,3.75) rectangle (2.52,4.25) node[midway,white] {TopAppBar};
        \foreach \y in {3.2,2.65,2.1,1.55}{\draw[petroleo!50,rounded corners=1pt,fill=petroleo!6] (0.2,\y) rectangle (2.4,\y+0.42);}
        \node[rotulo] at (1.3,1.25) {conteúdo};
        \fill[ambar,rounded corners=3pt] (1.85,0.72) rectangle (2.4,1.2) node[midway,white,font=\small] {+};
        \fill[petroleo!15] (0.08,0.1) rectangle (2.52,0.6) node[midway,petroleo] {NavigationBar};
      \end{tikzpicture}
      \begin{atencao}[\cod{innerPadding}]
        Aplique o padding recebido no conteúdo; senão ele fica escondido atrás das barras.
      \end{atencao}
    \end{column}
  \end{columns}
\end{frame}

\begin{frame}[fragile]{Tema Material 3}
  \begin{columns}[T]
    \begin{column}{0.53\textwidth}
\begin{codigo}[title=ui/theme/Theme.kt]{Kotlin}
private val Claro = lightColorScheme(
  primary = Color(0xFF0D3052),    // azul-petróleo
  secondary = Color(0xFFDE8A1C))  // âmbar
private val Escuro = darkColorScheme(
  primary = Color(0xFF9CC4EA),
  secondary = Color(0xFFE69A1F))
@Composable fun LabMonitorTheme(
  escuro: Boolean = isSystemInDarkTheme(),
  dinamica: Boolean = true,
  content: @Composable () -> Unit
) {
  val ctx = LocalContext.current
  val cores = when {
    dinamica && Build.VERSION.SDK_INT >= 31 ->
      if (escuro) dynamicDarkColorScheme(ctx)
      else dynamicLightColorScheme(ctx)
    escuro -> Escuro
    else -> Claro
  }
  MaterialTheme(cores, content = content)
}
\end{codigo}
    \end{column}
    \begin{column}{0.43\textwidth}
      \begin{conceito}[Três pilares]
        \textbf{Cores} (\cod{colorScheme}), \textbf{tipografia} (\cod{typography}) e \textbf{formas} (\cod{shapes}). Os componentes leem tudo do tema.
      \end{conceito}
      \begin{center}
      \begin{tikzpicture}
        \foreach \c/\n [count=\i] in {petroleo/primary,ambar/secondary,petroleoesc/tertiary,cinzaclaro/surface}{
          \fill[\c,rounded corners=2pt] (\i*1.05-1.05,0) rectangle ++(0.9,0.6);
          \node[rotulo] at (\i*1.05-0.6,-0.2) {\n};}
      \end{tikzpicture}
      \end{center}
      \begin{dica}[Use o tema, não cores fixas]
        \cod{MaterialTheme.colorScheme.primary} se adapta ao tema escuro e à cor dinâmica do Android 12+; \cod{Color.Red} não. O \emph{Material Theme Builder} gera as cores a partir de uma cor base.
      \end{dica}
    \end{column}
  \end{columns}
\end{frame}

\begin{frame}{Unidades, recursos e telas adaptáveis}
  \begin{columns}[T]
    \begin{column}{0.48\textwidth}
      \begin{conceito}[\cod{dp} e \cod{sp}]
        \begin{itemize}
          \item \cod{dp} (\emph{density-independent pixel}): tamanhos e espaços. 1\,dp $\approx$ 1/160 polegada.
          \item \cod{sp} (\emph{scale-independent pixel}): fontes; respeita o tamanho de fonte escolhido pelo usuário.
        \end{itemize}
      \end{conceito}
      \begin{atencao}[Alvo de toque]
        Elementos clicáveis precisam de pelo menos \textbf{48\,dp $\times$ 48\,dp}.
      \end{atencao}
    \end{column}
    \begin{column}{0.50\textwidth}
      \begin{conceito}[Telas adaptáveis]
        Use \termo{window size classes} para escolher o layout pelo espaço disponível, não pelo modelo do aparelho:
        \begin{itemize}
          \item \textbf{Compacta} (celular em retrato): uma coluna e \cod{NavigationBar}.
          \item \textbf{Média} (dobrável, tablet em retrato): \cod{NavigationRail}.
          \item \textbf{Expandida} (tablet, desktop): lista e detalhe lado a lado.
        \end{itemize}
      \end{conceito}
      \begin{dica}
        \cod{NavigationSuiteScaffold} troca automaticamente a barra inferior pela lateral conforme a largura.
      \end{dica}
    \end{column}
  \end{columns}
\end{frame}

% =============================================================================
\section{Fechamento}

\begin{frame}{Resumo}
  \begin{resumo}
    \begin{itemize}
      \item Em Compose, a UI é função do estado: descreva a tela e deixe o framework atualizá-la.
      \item Funções \cod{@Composable} são rápidas, sem efeitos colaterais e recebem um \cod{modifier}.
      \item \cod{Column}, \cod{Row} e \cod{Box} com \cod{Arrangement}, \cod{Alignment} e \cod{weight} resolvem quase todo layout.
      \item A ordem dos modifiers define o resultado.
      \item \cod{LazyColumn} com \cod{key} para listas; \cod{Scaffold} para a estrutura da tela.
      \item Cores, tipografia e formas vêm do \cod{MaterialTheme}: nada de cores fixas.
    \end{itemize}
  \end{resumo}
\end{frame}

\begin{frame}{Glossário (1/2)}
  \begin{glossario}
    UI declarativa & a interface é descrita como função do estado; o framework aplica as mudanças \\
    Composição & montar interfaces combinando componentes pequenos e independentes \\
    \cod{@Composable} & função que emite elementos de interface em vez de devolver um valor \\
    Recomposição & nova execução das funções que leram um estado alterado \\
    \emph{Skipping} & pular a função cujos parâmetros não mudaram desde a última execução \\
    Restrições & limites mínimos e máximos de largura e altura passados do pai ao filho \\
    \cod{Modifier} & cadeia ordenada que altera tamanho, espaço, desenho e interação \\
    \cod{LazyColumn} & lista que compõe apenas os itens visíveis \\
    Chave (\emph{key}) & identidade estável de um item entre recomposições \\
  \end{glossario}
\end{frame}

\begin{frame}{Glossário (2/2)}
  \begin{glossario}
    \cod{Scaffold} & estrutura padrão de uma tela Material: barras, FAB e conteúdo \\
    \emph{Token} de design & decisão visual nomeada pelo papel (\cod{primary}, \cod{titleMedium}) \\
    \cod{CompositionLocal} & valor disponível implicitamente para toda uma subárvore \\
    \cod{dp} e \cod{sp} & unidades independentes de densidade; o \cod{sp} segue a escala de fonte \\
    \emph{Window size class} & classe de largura da janela: compacta, média ou expandida \\
    \emph{Jank} & quadro perdido por exceder o orçamento de tempo \\
    Árvore de semântica & descrição da tela usada por leitores de tela como o TalkBack \\
  \end{glossario}
\end{frame}

\begin{frame}{Questões de revisão}
  \begin{questoes}
    \item Qual a diferença entre uma interface imperativa e uma declarativa? Por que a declarativa evita telas inconsistentes?
    \item Explique a regra ``restrições descem, tamanhos sobem, o pai posiciona'' com um exemplo numérico.
    \item Quais são as fases de um quadro no Compose e o que cada uma decide?
    \item O que é recomposição e quando uma função \cod{@Composable} é pulada?
    \item Por que uma função \cod{@Composable} não deve iniciar uma chamada de rede?
    \item Por que \cod{Modifier.padding(16.dp).background(c)} é diferente de \cod{Modifier.background(c).padding(16.dp)}?
    \item O que acontece em uma \cod{LazyColumn} sem \cod{key} quando o primeiro item, que estava expandido, é removido?
    \item Por que usar \cod{MaterialTheme.colorScheme.primary} em vez de \cod{Color(0xFF0D3052)} dentro de um componente?
  \end{questoes}
\end{frame}

\begin{frame}{Exercícios}
  \begin{exercicio}[Exercício 1 --- Layout]
    Reproduza um cartão de ``medidor'' com título, valor grande, unidade menor ao lado e uma barra de progresso (\cod{LinearProgressIndicator}) colorida conforme a faixa.
  \end{exercicio}
  \begin{exercicio}[Exercício 2 --- Modifiers]
    Crie um botão circular de 56\,dp com borda âmbar de 2\,dp, fundo azul-petróleo e ícone central. Troque a ordem de \cod{clip} e \cod{background} e explique o resultado.
  \end{exercicio}
  \begin{exercicio}[Exercício 3 --- Previews]
    Escreva \cod{@Preview}s do seu cartão nos temas claro e escuro, com fonte 1,5$\times$ e em largura de tablet.
  \end{exercicio}
  \begin{projeto}[LabMonitor --- Etapa 5]
    Crie a tela \textbf{Equipamentos}: \cod{Scaffold} com \cod{TopAppBar}, FAB e \cod{LazyColumn} de \cod{ItemEquipamento} usando os dados de \cod{DadosExemplo}. Aplique o tema com as cores do curso.
  \end{projeto}
\end{frame}

\begin{frame}{Para aprofundar}
  \small
  \begin{itemize}
    \item Tutorial oficial do Compose: \url{https://developer.android.com/develop/ui/compose/tutorial}
    \item Pensando em Compose: \url{https://developer.android.com/develop/ui/compose/mental-model}
    \item Layouts básicos: \url{https://developer.android.com/develop/ui/compose/layouts/basics}
    \item Modifiers: \url{https://developer.android.com/develop/ui/compose/modifiers}
    \item Material 3 em Compose: \url{https://developer.android.com/develop/ui/compose/designsystems/material3}
    \item Material Design 3: \url{https://m3.material.io}
  \end{itemize}
\end{frame}

\begin{frame}{Respostas das questões de revisão}
  \begin{respostas}
    \item No imperativo, o código altera componentes existentes e precisa sincronizá-los a cada evento; no declarativo, descreve-se a tela para cada estado e o framework aplica as diferenças. Há um só caminho do estado à tela, então nada fica esquecido.
    \item Na tela de 360\,dp, a \cod{Column} passa ao cartão largura de 0 a 360\,dp e altura livre; o cartão escolhe 360 $\times$ 96\,dp; a coluna o põe em $y = 0$ e o próximo em $y = 96$.
    \item Composição (o que mostrar), layout (medir e posicionar) e desenho (pintar); a GPU então gera os pixels. Uma mudança só de cor pode pular as duas primeiras fases.
    \item É a nova execução das funções que leram um estado alterado. A função é pulada quando seus parâmetros são estáveis e iguais aos da execução anterior.
    \item Porque ela pode rodar muitas vezes e em qualquer ordem: a chamada se repetiria a cada recomposição. Efeitos ficam em \cod{LaunchedEffect} (Módulo 07) ou no \cod{ViewModel}.
    \item Os modifiers se aplicam em cadeia, de fora para dentro: no primeiro, o espaçamento vem antes e o fundo cobre só o conteúdo; no segundo, o fundo cobre também o espaçamento.
    \item A correspondência é feita pela posição: o estado ``expandido'' passa para o item que ocupou a posição 0. Com \cod{key = \{ it.id \}}, o estado sai junto com o item removido.
    \item O \emph{token} se adapta ao tema escuro, à cor dinâmica e a mudanças da identidade visual; a cor fixa, não.
  \end{respostas}
\end{frame}

\encerramento{Módulo 06 --- Kotlin assíncrono: coroutines, dispatchers, Flow e StateFlow.}

\end{document}
